\section{Analysis}
This chapter analyses what the system must do and how its parts and data relate: the requirements,
the overall architecture, data flow diagrams and UML diagrams.

\subsection{Requirements Analysis}
The functional requirements were derived from the problems in Section~1.1 and from discussions with
students, faculty and the guide. Table~\ref{tab:fr} lists them with their acceptance criteria and
Table~\ref{tab:nfr} the non-functional targets.
\begingroup\small
\begin{longtable}{@{}p{1.1cm}p{3.4cm}p{6.2cm}p{4.9cm}@{}}
\caption{Functional requirements}\label{tab:fr}\\
\toprule
\textbf{ID} & \textbf{Requirement} & \textbf{Input $\rightarrow$ processing $\rightarrow$ output} & \textbf{Acceptance criterion}\\
\midrule
\endfirsthead
\toprule
\textbf{ID} & \textbf{Requirement} & \textbf{Input $\rightarrow$ processing $\rightarrow$ output} & \textbf{Acceptance criterion}\\
\midrule
\endhead
FR-01 & Sign-in with college e-mail and password & E-mail, password $\rightarrow$ Argon2id verify, account-status check, rate limit $\rightarrow$ 15-min access JWT + rotating refresh token & Wrong e-mail or password gives the same generic error; suspended account refused\\
FR-02 & Five roles enforced on the server & Request + token $\rightarrow$ roles loaded from \texttt{USER\_ROLE} on every request $\rightarrow$ allowed / 403 & No endpoint trusts a client-sent role; role change invalidates old tokens\\
FR-03 & Upload a timetable (PDF, image, XLSX/XLS, CSV, DOCX/DOC) & File $\rightarrow$ extension allow-list, content-signature check, size/page limits, private storage, SHA-256, job queued $\rightarrow$ archive record (\textsc{Queued}) & Upload never changes the primary timetable; retries with the same request id are idempotent\\
FR-04 & Extract class entries & Stored file $\rightarrow$ worker: page inspection, table/OCR extraction, sections, time rows, cell parsing, legend expansion, mapping $\rightarrow$ entries with provenance and status & Real college PDF: $\geq 95\%$ of class entries \textsc{Verified}; nothing guessed silently\\
FR-05 & Validate and classify the timetable & Entries $\rightarrow$ checks (missing fields, duplicates, time consistency, overlaps) $\rightarrow$ \textsc{Ready} / \textsc{Needs Review} / \textsc{Unusable} / \textsc{Failed} & Unusable timetables cannot become primary\\
FR-06 & Review and correct entries & Admin edits $\rightarrow$ before/after snapshot in \texttt{ENTRY\_CORRECTION}, audit row $\rightarrow$ corrected entry & Original value always recoverable; entry becomes \textsc{Verified} only by explicit reviewer action\\
FR-07 & Activate the official timetable & Admin action $\rightarrow$ settings row locked, primary switched, version incremented, users notified $\rightarrow$ new primary & Concurrent activations serialise; next request sees the new primary\\
FR-08 & Archives and selection & User picks an older version $\rightarrow$ selection remembered per domain $\rightarrow$ searches run on that archive with a banner & Archive browsing never mixes domains\\
FR-09 & Role dashboard & Login $\rightarrow$ current class, next class with countdown, free periods, later today $\rightarrow$ cards & Dashboard and typed questions give identical answers\\
FR-10 & Weekly timetable view & Day selection $\rightarrow$ entries for the user's batch and lab group (and parent division) $\rightarrow$ list with room/teacher & Lab group A1 does not see A2's labs\\
FR-11 & Natural-language search & English text $\rightarrow$ parse, build a query plan (operation, filters, lab/lecture, date/week, time window, sort), resolve entities in this timetable, validate, compile with role scope $\rightarrow$ result envelope & No SQL generated from text; counts, first/last and day-wise questions answered; ambiguous a.m./p.m., date or entity asks for clarification\\
FR-12 & Free-room and professor-availability search & Time/date, optional floor $\rightarrow$ room inventory minus scheduled rooms $\rightarrow$ free rooms / free periods & Uncertain entries treated as busy; incomplete inventory reported\\
FR-13 & Conflict detection & Timetable $\rightarrow$ interval overlap per professor, room, batch, excluding parallel lab groups and electives $\rightarrow$ conflict list & Real PDF: exactly the 2 genuine clashes reported\\
FR-14 & Personal timetables & Own upload $\rightarrow$ same pipeline, stored in personal tables $\rightarrow$ private timetable & Other users (including Admin) get 404 for it\\
FR-15 & Academic promotion & Admin preview + confirm $\rightarrow$ transactional update of placements with history $\rightarrow$ students moved to next year/batch & All-or-nothing; history preserved; concurrent promotions cannot double-apply\\
FR-16 & Master data management & Admin forms $\rightarrow$ validated CRUD of departments, years, batches, courses, rooms, aliases $\rightarrow$ records & Duplicates and invalid references rejected\\
FR-17 & User and account management & Admin $\rightarrow$ register student/staff, roles, suspend, issue access code $\rightarrow$ accounts & Suspension revokes sessions immediately\\
FR-18 & Notifications & Activation and job completion $\rightarrow$ in-app notification (push in store builds) $\rightarrow$ notification list & Users see only their own notifications\\
FR-19 & Audit log & Sensitive actions $\rightarrow$ append-only audit events $\rightarrow$ Admin audit screen & Uploads, activations, corrections, promotions, access denials logged\\
\bottomrule
\end{longtable}
\endgroup

\begin{table}[H]
\centering
\caption{Non-functional requirements}
\label{tab:nfr}
\small
\begin{tabularx}{\textwidth}{@{}p{1.3cm}p{2.7cm}X@{}}
\toprule
\textbf{ID} & \textbf{Category} & \textbf{Requirement}\\
\midrule
NFR-01 & Performance & Median search response $<300$~ms for one user; extraction of a 10-page PDF $<10$~s.\\
NFR-02 & Scalability & Stateless API that scales horizontally by adding processes; workers scale independently.\\
NFR-03 & Availability, reliability & Correct answers under concurrent use; idempotent uploads; failed jobs retried up to 3 times; health endpoint for the platform.\\
NFR-04 & Security & Argon2id passwords, JWT + refresh-token rotation and reuse detection, RBAC, IDOR protection, rate limiting, parameterised SQL only, secrets never in the app bundle.\\
NFR-05 & Privacy & Personal timetables visible only to their owner; no free text in logs.\\
NFR-06 & Usability, accessibility & Readable in light and dark mode; large tap targets; plain-English questions; clarifications with tappable choices.\\
NFR-07 & Maintainability & Modular layers (API, services, query, extraction); typed code (Pydantic, TypeScript strict); automated tests; documented decisions.\\
NFR-08 & Portability & One code base for Android, iOS and web; Docker image for any Linux host.\\
NFR-09 & Resource utilisation & One API worker $<300$~MB RAM; database $<50$~MB for one institute's timetables.\\
NFR-10 & Auditability & Every entry keeps raw text, page, parser version and status; every correction keeps before/after.\\
\bottomrule
\end{tabularx}
\end{table}


\subsection{System Architecture}
Figure~\ref{fig:arch} shows the architecture. One React Native app serves Android, iOS and the web.
It talks over HTTPS to a stateless FastAPI back end that authenticates every request, loads the
user's roles from the database and resolves the search context on the server. Uploads are queued in
PostgreSQL and processed by a background worker, so the API is never blocked by OCR. PostgreSQL is the
single source of truth.
\begin{figure}[H]
\centering
\begin{tikzpicture}[node distance=4mm and 6mm]
% clients
\node[ext,minimum width=3.3cm] (android) {Android app};
\node[ext,minimum width=3.3cm,right=of android] (ios) {iPhone app};
\node[ext,minimum width=3.3cm,right=of ios] (web) {Laptop browser\\(web build)};
\node[draw,dashed,rounded corners,fit=(android)(ios)(web),inner sep=6pt,
      label={[font=\scriptsize\bfseries]above:React Native / Expo client (one TypeScript code base)}] (cl) {};
% api
\node[blk,below=13mm of ios,minimum width=12.6cm,minimum height=9mm] (mw) {Middleware: request id, JSON logs, CORS, error envelope, rate limiting};
\node[blk,below=3mm of mw.south west,anchor=north west,minimum width=2.9cm,minimum height=13mm] (sec) {Security\\JWT $\cdot$ RBAC\\policy};
\node[blk,right=2mm of sec,minimum width=2.9cm,minimum height=13mm] (routers) {REST routers\\auth, me, timetables,\\search, admin};
\node[blk,right=2mm of routers,minimum width=3.2cm,minimum height=13mm] (svc) {Services\\context, uploads,\\conflicts, academics};
\node[blk,right=2mm of svc,minimum width=3cm,minimum height=13mm] (qe) {Query planner\\NLP $\rightarrow$ query plan\\$\rightarrow$ QT01--QT23};
\node[draw,thick,rounded corners,fit=(mw)(sec)(qe),inner sep=5pt,
      label={[font=\scriptsize\bfseries,anchor=south west]north west:FastAPI application (stateless, N processes)}] (api) {};
% bottom
\node[store,below=16mm of routers,minimum width=3.6cm,minimum height=11mm,font=\small] (db) {PostgreSQL 16\\31 tables};
\node[blk,below=16mm of qe,minimum width=4cm,minimum height=11mm] (worker) {Background worker\\Extraction pipeline\\(PyMuPDF, Tesseract, Office)};
\node[store,left=10mm of db,minimum width=3cm,minimum height=11mm,font=\small] (files) {Private file\\storage};
\draw[darr] (cl.south) -- node[lbl,right]{HTTPS, JSON (REST /api/v1)} (mw.north -| cl.south);
\draw[darr] (api.south -| db.north) -- node[lbl,right]{SQL (bound params)} (db.north);
\draw[arr] (worker.west) -- node[lbl,above]{claims jobs} node[lbl,below,font=\tiny]{FOR UPDATE SKIP LOCKED} (db.east);
\draw[arr] (api.south -| files.north) -- node[lbl,left]{originals} (files.north);
\draw[arr] (files.south) to[bend right=25] node[lbl,below]{read file} (worker.south);
\end{tikzpicture}
\caption{Overall system architecture (block diagram)}
\label{fig:arch}
\end{figure}


\subsection{Data Flow Diagrams}
The context diagram (Figure~\ref{fig:dfd0}) shows the system and its external entities. Level~1
(Figure~\ref{fig:dfd1}) splits it into the main processes: authentication, upload and extraction,
review and activation, conflict detection and question answering, all sharing the timetable store.
\begin{figure}[H]
\centering
\begin{adjustbox}{max width=\textwidth}
\begin{tikzpicture}
\node[proc,minimum size=34mm,font=\small] (sys) {0.0\\Dynamic\\Timetable\\System};
\node[ext,above left=8mm and 26mm of sys] (stu) {Student};
\node[ext,above right=8mm and 26mm of sys] (prof) {Professor};
\node[ext,below left=8mm and 26mm of sys] (adm) {Admin};
\node[ext,below right=8mm and 26mm of sys] (hod) {HOD /\\Principal};
\draw[arr] (stu.east) to[bend left=10] node[lbl,above,sloped]{questions, personal file} (sys.north west);
\draw[arr] (sys.west) to[bend left=10] node[lbl,below,sloped]{schedule, answers} (stu.south);
\draw[arr] (prof.west) to[bend right=10] node[lbl,above,sloped]{questions} (sys.north east);
\draw[arr] (sys.east) to[bend right=10] node[lbl,below,sloped]{schedule, free rooms} (prof.south);
\draw[arr] (adm.east) to[bend right=10] node[lbl,below,sloped]{timetable file, corrections} (sys.south west);
\draw[arr] (sys.west) to[bend right=30] node[lbl,above,sloped]{validation summary} (adm.north);
\draw[arr] (sys.south east) to[bend left=10] node[lbl,below,sloped]{conflicts, room usage} (hod.west);
\end{tikzpicture}
\end{adjustbox}
\caption{DFD level 0 (context diagram)}
\label{fig:dfd0}
\end{figure}

\begin{figure}[H]
\centering
\begin{adjustbox}{max width=\textwidth}
\begin{tikzpicture}
\node[ext] (adm) at (0,0) {Admin};
\node[proc] (p1) at (3.4,0) {1.0\\Receive \&\\archive};
\node[proc] (p2) at (6.9,0) {2.0\\Extract \&\\validate};
\node[proc] (p3) at (10.6,0) {3.0\\Review \&\\activate};
\node[store] (d1) at (5.15,2.3) {D1 Uploaded files};
\node[store] (d3) at (0,-3.4) {D3 Users, roles,\\preferences};
\node[proc] (p4) at (3.4,-3.4) {4.0\\Resolve\\context};
\node[proc] (p5) at (6.9,-3.4) {5.0\\Answer query\\(query plan\\$\rightarrow$ SQL)};
\node[store] (d2) at (10.6,-3.4) {D2 Timetables\\\& entries};
\node[proc] (p6) at (10.6,-6.6) {6.0\\Detect\\conflicts};
\node[ext] (hod) at (14.4,-6.6) {HOD /\\Principal};
\node[ext] (usr) at (6.9,-6.9) {Student /\\Professor};
\draw[arr] (adm) -- node[lbl,above]{file} (p1);
\draw[arr] (p1) -- node[lbl,above]{job} (p2);
\draw[arr] (p1) |- node[lbl,pos=0.75,above]{original} (d1.west);
\draw[arr] (d1.east) -| (p2);
\draw[arr] (p2) -- node[lbl,above]{entries} (p3);
\draw[arr] (adm.north) -- ++(0,3.4) -| node[lbl,pos=0.25,above]{corrections, activation} (p3.north);
\draw[arr] (p3) -- node[lbl,right]{verified data} (d2);
\draw[arr] (d3) -- (p4);
\draw[arr] (p4) -- node[lbl,above]{context} (p5);
\draw[arr] (d2) -- node[lbl,above]{entries} (p5);
\draw[arr] (d2) -- node[lbl,right]{entries} (p6);
\draw[arr] (p6) -- node[lbl,above]{clashes} (hod);
\draw[arr] ([xshift=-5mm]usr.north) -- node[lbl,left]{question} ([xshift=-5mm]p5.south);
\draw[arr] ([xshift=5mm]p5.south) -- node[lbl,right]{answer} ([xshift=5mm]usr.north);
\end{tikzpicture}
\end{adjustbox}
\caption{DFD level 1}
\label{fig:dfd1}
\end{figure}


\subsection{Use Case Diagram}
Figure~\ref{fig:usecase} shows what each role can do. Every logged-in user can view any official
timetable and ask questions; personal timetables belong only to their owner; clash reports are for
the HOD, Principal and Admin; upload, correction and activation are Admin-only.
\begin{figure}[H]
\centering
\begin{adjustbox}{max width=\textwidth}
\begin{tikzpicture}[uc/.style={draw,ellipse,thick,fill=boxfill,font=\scriptsize,align=center,minimum width=33mm,minimum height=7mm,inner sep=1pt},
  actor/.style={font=\small,align=center}]
\node[draw,thick,rounded corners,minimum width=8.6cm,minimum height=12.3cm,label={[font=\small\bfseries]above:Dynamic Timetable}] (b) at (0,0) {};
\foreach \n/\t/\y in {u1/Sign in/5.4, u2/View dashboard \& week/4.5, u3/Ask question (NL search)/3.6, u4/Find free rooms/2.7,
  u5/Manage personal timetable/1.8, u6/View own teaching \& free time/0.9, u7/Review conflicts/0, u8/View floor activity/-0.9,
  u9/Upload official timetable/-1.8, u10/Review \& correct entries/-2.7, u11/Activate timetable/-3.6, u12/Manage master data \& users/-4.5, u13/Promote students/-5.4}
  { \node[uc] (\n) at (0,\y) {\t}; }
\node[actor] (s) at (-6.2,4) {\stick\\\textbf{Student}};
\node[actor] (p) at (-6.2,0.3) {\stick\\\textbf{Professor}};
\node[actor] (h) at (6.2,0.6) {\stick\\\textbf{HOD /}\\\textbf{Principal}};
\node[actor] (a) at (6.2,-3.6) {\stick\\\textbf{Admin}};
\foreach \u in {u1,u2,u3,u4,u5} \draw (s.east) -- (\u.west);
\foreach \u in {u1,u2,u3,u4,u6} \draw (p.east) -- (\u.west);
\foreach \u in {u1,u2,u7,u8} \draw (h.west) -- (\u.east);
\foreach \u in {u1,u9,u10,u11,u12,u13} \draw (a.west) -- (\u.east);
\end{tikzpicture}
\end{adjustbox}
\caption{Use case diagram}
\label{fig:usecase}
\end{figure}


\subsection{Class Diagram}
Figure~\ref{fig:class} shows the core domain: a timetable has many entries, and each entry refers to a
course, a teacher (staff), a room and a batch; batches form a hierarchy (division $\rightarrow$ lab
batch).
\begin{figure}[H]
\centering
\resizebox{\textwidth}{!}{%
\begin{tikzpicture}[cls/.style={draw,thick,rectangle split,rectangle split parts=3,fill=boxfill,
  font=\scriptsize,align=left,text width=40mm,rectangle split part fill={spitblue!25,boxfill,white}}]
\node[cls] (ua) at (0,0) {\textbf{UserAccount}
\nodepart{second}user\_id: UUID\\college\_email: str\\password\_hash: str\\account\_status\\token\_version: int
\nodepart{third}+roles(): set[Role]};
\node[cls] (st) at (-5,-4.2) {\textbf{Student}
\nodepart{second}uid: str\\student\_status
\nodepart{third}+current\_placement()};
\node[cls] (sf) at (0,-4.2) {\textbf{Staff}
\nodepart{second}short\_code: str\\designation
\nodepart{third}};
\node[cls] (it) at (5.4,0) {\textbf{InstitutionalTimetable}
\nodepart{second}timetable\_id: UUID\\processing\_status\\parser\_version\\effective\_from: date\\data\_version: int
\nodepart{third}+activate()\\+conflicts(): list};
\node[cls] (ie) at (5.4,-4.6) {\textbf{InstitutionalEntry}
\nodepart{second}day\_of\_week, start/end\\course/staff/room/batch\\raw\_extracted\_text\\source\_page\\verification\_status
\nodepart{third}+correct(change)};
\node[cls] (job) at (10.6,0) {\textbf{ProcessingJob}
\nodepart{second}status, attempts\\heartbeat\_at
\nodepart{third}+claim()\\+fail(code)};
\node[cls] (qe) at (-5,0) {\textbf{QueryEngine}
\nodepart{second}registry: Q01..Q30\\templates: QT01..QT22
\nodepart{third}+search(ctx, text)\\-classify(), -resolve()};
\node[cls] (b) at (10.6,-4.6) {\textbf{Batch}
\nodepart{second}code (SE-A1)\\parent\_batch\_id
\nodepart{third}+ancestors()};
\draw[thick] (ua.south) -- node[lbl,pos=.85,right]{0..1} (sf.north);
\draw[thick] (ua.south west) -- node[lbl,pos=.85,left]{0..1} (st.north);
\draw[thick,-{Diamond[open,length=3mm]}] (ie.north) -- node[lbl,pos=0.2,right]{1..*} (it.south);
\draw[thick] (job.west) -- node[lbl,above]{1} (it.east);
\draw[thick,dashed,-{Stealth}] (qe.east) -- node[lbl,above]{reads} (ua.west);
\draw[thick] (ie.east) -- node[lbl,above]{*\quad 0..1} (b.west);
\draw[thick] (sf.east) -- node[lbl,above]{0..1\quad *} (ie.west);
\draw[thick] (b.north) .. controls +(0.8,1) and +(1.6,0.2) .. node[lbl,right]{parent} (b.east);
\end{tikzpicture}}
\caption{Simplified class diagram of the core domain}
\label{fig:class}
\end{figure}


\subsection{Sequence Diagram}
Figure~\ref{fig:seq} follows one question from the app to the database and back: authentication,
context resolution, planning, validation and execution of a parameterised query.
\begin{figure}[H]
\centering
\begin{tikzpicture}[
  ll/.style={draw,thick,fill=boxfill,font=\scriptsize,minimum width=19mm,minimum height=6mm},
  msg/.style={-{Stealth},thick,font=\scriptsize},
  ret/.style={-{Stealth},dashed,font=\scriptsize}]
\foreach \x/\n/\t in {0/c/App, 2.6/a/API router, 5.2/s/Security, 7.8/x/Context, 10.4/q/Query engine, 13.0/d/PostgreSQL}{
  \node[ll] (\n) at (\x,0) {\t};
  \draw[dashed,gray] (\x,-0.3) -- (\x,-9.3);}
\draw[msg] (0,-0.9) -- node[above]{POST /search \{query\}} (2.6,-0.9);
\draw[msg] (2.6,-1.6) -- node[above]{verify JWT} (5.2,-1.6);
\draw[msg] (5.2,-2.2) -- node[above]{load user+roles} (13,-2.2);
\draw[ret] (13,-2.7) -- node[above]{principal} (5.2,-2.7);
\draw[msg] (2.6,-3.4) -- node[above]{resolve(domain, timetable, tz)} (7.8,-3.4);
\draw[msg] (7.8,-4.0) -- node[above]{preferences, settings} (13,-4.0);
\draw[ret] (7.8,-4.6) -- node[above]{context} (2.6,-4.6);
\draw[msg] (2.6,-5.3) -- node[above]{search(ctx, text)} (10.4,-5.3);
\draw[thick] (10.4,-5.6) -- ++(0.6,0) -- node[right,font=\scriptsize,align=left]{parse, classify,\\resolve entities,\\validate params} ++(0,-0.9) -- ++(-0.6,0);
\draw[msg] (10.4,-7.0) -- node[above]{template QTnn (bound params)} (13,-7.0);
\draw[ret] (13,-7.6) -- node[above]{rows} (10.4,-7.6);
\draw[ret] (10.4,-8.2) -- node[above]{envelope: status, results, warnings} (2.6,-8.2);
\draw[ret] (2.6,-8.9) -- node[above]{200 / 422 clarification} (0,-8.9);
\end{tikzpicture}
\caption{Sequence diagram for a natural-language search}
\label{fig:seq}
\end{figure}


\subsection{ER Diagram}
Figure~\ref{fig:er} shows the main entities of the database. The schema has 31 tables in third normal
form; course, staff, room and batch details are stored once and referenced by foreign keys.
\begin{figure}[H]
\centering
\begin{adjustbox}{max width=\textwidth}
\begin{tikzpicture}[ent/.style={draw=spitblue,thick,fill=boxfill,font=\scriptsize\bfseries,minimum width=27mm,minimum height=7mm,align=center},
  rel/.style={draw,thick,diamond,aspect=2,fill=yellow!15,font=\tiny,inner sep=1pt,align=center}]
\node[ent] (ua) at (0,0) {USER\_ACCOUNT};
\node[ent,minimum width=16mm] (role) at (-6.2,0) {ROLE};
\node[rel] (has) at (-3.4,0) {USER\_ROLE};
\node[ent] (stu) at (-2.5,-2.2) {STUDENT};
\node[ent] (stf) at (2.5,-2.2) {STAFF};
\node[ent] (sah) at (-2.5,-4.4) {STUDENT\_ACADEMIC\\\_HISTORY};
\node[ent] (batch) at (-2.5,-6.8) {BATCH};
\node[ent] (dept) at (-6.6,-5.6) {DEPARTMENT};
\node[ent] (ay) at (-6.6,-3.4) {ACADEMIC\_YEAR};
\node[ent] (itt) at (5.2,-4.4) {INSTITUTIONAL\\\_TIMETABLE};
\node[ent] (ite) at (2.5,-6.8) {INSTITUTIONAL\\\_TIMETABLE\_ENTRY};
\node[ent] (room) at (2.5,-9) {ROOM};
\node[ent] (course) at (-2.5,-9) {COURSE};
\node[ent] (set) at (8.6,-4.4) {SETTINGS\\(single row)};
\node[ent] (ptt) at (5.2,0) {PERSONAL\\\_TIMETABLE};
\node[ent] (pte) at (8.6,0) {PERSONAL\\\_ENTRY};
\node[ent] (job) at (8.6,-6.8) {PROCESSING\\\_JOB};
\node[ent] (cor) at (6,-9) {ENTRY\\\_CORRECTION};
\draw (role) -- node[lbl,above]{M} (has) -- node[lbl,above]{N} (ua);
\draw (ua) -- node[lbl,left]{1 : 0..1} (stu);
\draw (ua) -- node[lbl,right]{1 : 0..1} (stf);
\draw (stu) -- node[lbl,left]{1 : N} (sah);
\draw (sah) -- node[lbl,left]{N : 1} (batch);
\draw (sah) -- node[lbl,above,sloped]{N : 1} (ay);
\draw (batch) -- node[lbl,above,sloped]{N : 1} (dept);
\draw (batch) .. controls +(-1.6,-0.9) and +(-1.6,0.9) .. node[lbl,left]{parent} (batch);
\draw (itt) -- node[lbl,above,sloped]{1 : N} (ite);
\draw (ite) -- node[lbl,above]{N : 1} (batch);
\draw (ite) -- node[lbl,right]{N : 1} (room);
\draw (ite) -- node[lbl,above,sloped]{N : 1} (course);
\draw (ite) -- node[lbl,right]{N : 1} (stf);
\draw (set) -- node[lbl,above]{0..1} (itt);
\draw (ua) -- node[lbl,above]{1 : N owns} (ptt);
\draw (ptt) -- node[lbl,above]{1 : N} (pte);
\draw (job) -- node[lbl,above,sloped]{N : 1} (itt);
\draw (cor) -- node[lbl,above,sloped]{N : 1} (ite);
\end{tikzpicture}
\end{adjustbox}
\caption{ER diagram (main entities; personal sections, notifications, tokens and audit omitted)}
\label{fig:er}
\end{figure}

\newpage
